% Transcription — fonds Alexandre Grothendieck, shelfmark 67, batch 3 (pages 41-60).
% Established from the facsimile archives/batches/67/batch-03.pdf (a mirror of
% https://grothendieck.umontpellier.fr/67.pdf), per the skill
% transcribe-grothendieck. The batch file carries no cover sheet: PDF page N
% is page 40+N of the folder (the Montpellier footer prints « 41 » ... « 60 »).
%
% Pass: Opus 5.5 (claude-opus-5-5), 2026-09-22 — first pass, unchecked against the pages by a human.
% Revised 2026-09-23 (Opus 5.5), against the facsimile at 400 dpi: the
% circled numbers of pages 41-43 are « 16 », « 17 », « 18 », not « 10 »,
% « 11 », « 12 » — they continue batch 2's 1-15 — and those of pages 44-45
% are « 1′ » (probable) and « 2′ », not « 17 » and « 21 »; the claim that
% his sheets 13-16 and 18-20 are missing is withdrawn. The notes of pages
% 41-45 say what is on each leaf and what is inferred. The note of page 53
% no longer says pages 53-60 are numbered 1 to 5: his 1-4 stand on pages 53,
% 55, 57, 59 and the digit on page 60 is scratched out.
%
% Eighteen of the twenty pages are transcribed. Pages 46 and 51 are blank
% leaves and are absent. \page{} numbers the position in the folder, as
% batches 1 and 4 do: the archivists' pencil numbers bottom-left run three
% ahead here (page 44 is pencilled 47, page 53 is pencilled 56).
%
% What the batch is.
%   41-45  Ink on punched sheets. Pages 41-43, tractor-feed like batch 2's
%          run, carry his circled « 16 », « 17 », « 18 » top left: they
%          continue his sheets 1-15 (pages 26-40), and page 40's last
%          sentence runs on into page 41. Pages 44 and 45, on another paper,
%          carry circled numbers top centre, read « 1′ » and « 2′ » (the
%          second clear; on 44 the sign after the 1 is short and raised like
%          45's prime, « 17 » not wholly excluded): inserts. Where they go is
%          inferred from the text, not written: 44 after his sheet 1 (page
%          26), 45 after his sheet 2 (page 27), page 28 continuing 45's last
%          sentence. Content: the
%          distinguished (canonical) metric on a conformal surface with
%          boundary, seen as the quotient of X~ by an anti-involution σ with
%          fixed points. Antiinvolutions of C (z ↦ a z̄ + b, fixed line α.R,
%          the half-plane H̄ with the real affine group), of P^1_C (the disc);
%          the four exceptional cases a) a') b) b') ; Corollaire 3 (X ⊂ X'
%          0-connected is hyperbolic except for P^1, E^1, G_m with ≤ 2, 1, 0
%          points removed, or genus 1); non-orientable hyperbolic surfaces;
%          constant curvature -1; the elliptic case begins (page 45 breaks
%          off mid-sentence, on « La constante multiplicative »).
%   47-50  Computer-listing paper (printed 16 AUG 82 / 11-08-82), dated
%          « 15.1. »: Teichmüller isomorphisms X ≃_T Y between a smooth curve
%          of type (g,ν) over Q̄ and a stable curve with normal propagator N,
%          outer isomorphisms π_1(X) → π_1([Y,N]); « À examiner » a) b) c)
%          (anabelian conjecture for stable curves, recovering singular and
%          marked points from inertia classes, anabelian compactifications
%          after « MD » of M_{g,ν}); generators and relations for Teichmüller
%          groupoids: pants/triads, Lego boxes, involutions, with
%          SL(2,Z) = <ρ,σ | ρ^3 = σ^2, ρ^6 = 1> and GL(2,Z) by three
%          reflections. Page 50 is a scratch sheet written in every
%          direction (scanned upside down): exact sequences
%          1 → Π_{0,3} → T^+_{0,4} → S_4 → 1 etc., notation lists, sphere
%          sketches — transcribed as a list, layout not reproduced.
%   52     Scratch: toric circles ↔ oriented real planes ↔ unitary complex
%          lines; U-torsors; the datum ((D',D''),u) ↔ a duality T' ⊗ T'' plus
%          a disc. Heavily overwritten.
%   53-60  « 11.1.84 », « Chirurgie conforme », his pp. 1-4, written at the
%          head of the rectos 53, 55, 57, 59 (on 59 the « 4 » follows a
%          struck digit; 54, 56, 58 carry none); the digit at the head of
%          page 60 is scratched out and not read; batch 4 opens on his « 5 »
%          at page 61: (1)
%          conformal surfaces with boundary, Théorème 1 (the conformal
%          structure is determined by the interior; Lemme 1, Schwarz
%          reflection), [Thm 2: compact conformal surfaces with boundary ≃
%          smooth proper real algebraic curves]; (2) Découpage along a
%          rectifiable 1-complex K, Théorème 3, Lemme 2 (uniformising the two
%          banks of an arc), Scholie; (3) Recollement, Th. 4 (equivalence
%          (X,K,S) ↔ (X~, S~, σ)), Lemme 3 in a struck and a kept version;
%          (4) Recollements bord à bord de disques. The last sentence of page
%          60 runs on into batch 4.
%
% Notation fixed for the batch: H̄ (and, from page 55 by his own NB, H) the
% closed upper half-plane, \mathbb{H} where he writes it double-struck
% (pages 42-44); \mathfrak{X} for his script X (the complex double X~);
% \mathring{D} for the open unit disc; \mathcal{T} for his script T
% (Teichmüller groupoids), \mathbb{T} for the double-struck T (groups);
% \operatorname{Int} for Int.
%
% Apparatus: \ill{} and \uncertain{} counts are given in the closing report
% of the pass; margins written obliquely were turned upright before reading.
\documentclass[11pt,a4paper]{article}
% Le préambule commun à toutes les transcriptions du fonds Grothendieck.
%
% Il tient en une idée : **l'incertitude se compose, elle ne se résout pas.**
% Une transcription qui lisse un mot douteux a détruit la seule chose pour
% laquelle on la faisait — un lecteur doit pouvoir savoir, à chaque endroit,
% ce qui était sur le feuillet et ce qui a été ajouté. Les six macros qui
% suivent sont donc l'appareil critique, et non de la décoration.
%
% Les mêmes commandes sont reconnues par `scripts/render.mjs`, qui produit la
% vue de lecture du site. Une macro ajoutée ici doit l'être aussi là-bas,
% faute de quoi le rendu échouera bruyamment — ce qui est voulu : un rendu
% silencieusement faux est pire qu'un rendu absent.

\ProvidesPackage{grothendieck}[2026/08/08 Transcriptions du fonds Grothendieck]

\usepackage{amsmath,amssymb,amsthm}
\usepackage{mathrsfs}
\usepackage{stmaryrd}
% Les diagrammes commutatifs sont partout dans ce fonds : c'est la langue
% même de la géométrie algébrique de Grothendieck, et une transcription qui
% les rendrait en prose ne transcrirait rien.
\usepackage{tikz-cd}
% Français par défaut — c'est la langue des feuillets — mais l'anglais est
% chargé aussi : la traduction et le résumé sont des documents anglais, et la
% typographie française (espaces avant les deux-points) y serait une faute.
% Ces documents basculent par \selectlanguage{english} en tête de corps.
\usepackage[english,french]{babel}
\usepackage{fontspec}
\usepackage{geometry}
\geometry{a4paper,margin=25mm,marginparwidth=28mm}
\usepackage{marginnote}
\usepackage{xcolor}
% ulem, not soul. `soul`'s \st and \ul are built on a character-by-character
% scan that XeLaTeX's font machinery breaks: the document fails with an
% undefined control sequence rather than degrading. ulem does the same two jobs
% and is Unicode-engine safe. `normalem` stops it hijacking \emph, which the
% transcriptions use for Grothendieck's own emphasis.
\usepackage[normalem]{ulem}
\usepackage{hyperref}
\hypersetup{colorlinks=true,linkcolor=black,urlcolor=black,pdfborder={0 0 0}}

% --- Métadonnées du lot -----------------------------------------------------
% Reprises telles quelles par le script de rendu, qui les affiche en tête de
% la vue de lecture. Elles fixent ce que le fichier prétend couvrir : sans
% elles, une transcription flotte sans point d'attache dans un fonds de seize
% mille feuillets.

\newcommand{\folder}[1]{\gdef\grothFolder{#1}}
\newcommand{\batch}[1]{\gdef\grothBatch{#1}}
\newcommand{\leaves}[2]{\gdef\grothFirst{#1}\gdef\grothLast{#2}}
\newcommand{\foldertitle}[1]{\gdef\grothFoldertitle{#1}}
\gdef\grothFolder{?}\gdef\grothBatch{?}\gdef\grothFirst{?}\gdef\grothLast{?}\gdef\grothFoldertitle{}

% --- Le résumé --------------------------------------------------------------
%
% Il ouvre la lecture modernisée, sous le titre « Résumé », et il est composé
% en petit. Ce n'est pas un ornement : c'est le seul passage du document qui
% ne soit pas des mathématiques, et un lecteur doit pouvoir le reconnaître
% avant de l'avoir lu — puis le sauter s'il connaît déjà le sujet, ou s'y
% arrêter s'il ne le connaît pas. Le filet qui le suit marque la reprise du
% corps.

\newenvironment{resume}
  {\small\setlength{\parskip}{0.45em}}
  {\par\medskip\hrule\medskip}

% --- Mots-clés ---------------------------------------------------------------
%
% En anglais, en clôture du résumé de la lecture modernisée : le vocabulaire
% moderne sous lequel chercher ce que les feuillets construisent. Le manifeste
% du site (`npm run manifest`) les extrait pour étiqueter la cote entière —
% cette ligne est la source unique des tags, il n'y a pas de fichier à côté.
\newcommand{\keywords}[1]{\par\smallskip\noindent
  {\footnotesize\textsc{Keywords} --- \textit{#1}}\par}

% --- L'appareil critique ----------------------------------------------------

% Le numéro de feuillet, en marge. C'est l'ancre qui permet de régler un
% désaccord sur une formule en regardant *un* feuillet plutôt qu'une chemise
% de six cents. Le site s'en sert aussi pour tourner les pages du fac-similé
% quand on fait défiler la transcription.
\newcommand{\leaf}[1]{%
  \marginnote{\footnotesize\color{gray!70}\textsf{#1}}%
  \label{leaf:#1}\ignorespaces}

% Illisible : on ne devine pas. Un mot inventé qui se lit comme les autres est
% le pire résultat possible.
\newcommand{\ill}{\textcolor{red!60!black}{[\,\ldots\,]}}

% Lecture douteuse : le mot est donné, et signalé comme incertain.
\newcommand{\uncertain}[1]{\uline{#1}}

% Ajout de l'éditeur — une lettre manquante, un mot sous-entendu.
\newcommand{\add}[1]{\textcolor{blue!50!black}{[#1]}}

% Biffé par Grothendieck lui-même. Ce qu'il a rayé fait partie du document :
% une rature dit souvent où la pensée a changé de direction.
\newcommand{\struck}[1]{\sout{#1}}

% Note du transcripteur, sur l'état matériel du feuillet.
\newcommand{\note}[1]{\par\smallskip{\small\color{gray!60!black}\textit{[#1]}}\par\smallskip}

% Note marginale de Grothendieck, distincte du corps du texte.
\newcommand{\marginal}[1]{\par\smallskip\noindent\hspace{1em}%
  {\small\color{gray!60!black}#1}\par\smallskip}

% --- Titre ------------------------------------------------------------------

\AtBeginDocument{%
  \begin{center}
    {\large\bfseries Fonds Alexandre Grothendieck --- cote n\textsuperscript{o}~\grothFolder}\\[2pt]
    {\itshape\grothFoldertitle}\\[4pt]
    {\small feuillets \grothFirst--\grothLast\ (lot \grothBatch)}\\[2pt]
    {\footnotesize\color{gray!60!black}Transcription établie d'après le fac-similé numérisé
     par l'Université de Montpellier.\\
     Les lectures incertaines sont soulignées, les illisibles notées [\,\ldots\,],
     les ajouts entre crochets.}
  \end{center}
  \vspace{1em}}

\endinput


\folder{67}
\batch{3}
\pages{41}{60}
\dating{[à partir de 1977]-1984}
\watermark{Édition de démonstration}
\foldertitle{Chirurgie des surfaces conformes : notes manuscrites (s.d., 1983-1984), lettres (1984)}

\begin{document}

\page{41}
\note{en haut à gauche, un « 16 » cerclé, de sa main : il fait suite au « 15 » de la page 40 (lot 2), dont la dernière phrase se poursuit ici.}
\ill{} dans quelques cas que si on le complète par la \uncertain{présomption} que l'on peut \ill{} une métrique « distinguée » sur \uncertain{$X^{\sigma}$} qui soit stable par $\sigma$.

Considérons maintenant les seuls cas nouveaux --- ceux où $X$ est connexe et $\partial X \neq \emptyset$, compactes : un \uncertain{$X^{\sim}$} $= \mathfrak{X}$ connexe, muni d'une antiinvolution $\sigma$ telles que $\mathfrak{X}^{\sigma} \neq \emptyset$. Les seuls cas particuliers pour l'existence d'une métrique canonique sur $X$, sont les suivants :

a) $\mathfrak{X} \simeq \mathbb{C}$. \struck{\ill{}} OPS $\mathfrak{X} = \mathbb{C}$. Les anti-aut.\ de $\mathbb{C}$ sont
\[
  z \longmapsto a\bar{z} + b \qquad (a \in \mathbb{C}^{*},\ b \in \mathbb{C})
\]
Si on prend une antiinvolution à pts fixes, OPS que $0$ est pt fixe, donc $b = 0$, donc \struck{l'application $z \mapsto a\bar{z}$}
\[
  \sigma z = a\bar{z}
\]
avec $a\bar{a} = 1$ i.e.\ $|a| = 1$ i.e.\ $a \in \mathbb{U}$ (puisque $\sigma^{2} = \mathrm{id}$) donc \struck{l'ens.\ des pts fixes est} si $a = \alpha^{2}$, $\alpha \in \mathbb{U}$,
\[
  \sigma z = \alpha^{2}\bar{z} = \alpha/\bar{\alpha}\,\bar{z}
\]
et l'ens.\ des pts fixes est
\[
  \mathbb{C}^{\sigma} = \alpha.\mathbb{R}
\]
\marginal{symétrie par rapport à la droite $\alpha.\mathbb{R}$}

\page{42}
\note{en haut à gauche, un « 17 » cerclé, de sa main.}
En fait, les antiinvolutions de $\mathbb{C}$ ont toutes des pts fixes, et ce sont les symétries orth.\ p.r.\ aux droites réelles (\uncertain{qu'on suppose} passant par l'origine) dans $\mathbb{C}$. La surface à bord conforme correspondante est donc : un \emph{demi-plan} (avec bord). On note que les métriques sur $\mathfrak{X}$, quasi-invariantes par autom.\ conformes, et invariantes par autom.\ d'ordre fini, en particulier par $\sigma$. Donc il y a sur $X$ une \uncertain{donnée} métrique canonique, mod.\ ct.

NB \struck{\ill{}} Le modèle des $X$ envisagés est le demi-plan supérieur \struck{\ill{}}
\[
  \overline{\mathbb{H}} = \{z \in \mathbb{C} \mid \Im z \geq 0\}.
\]
Le groupe des autom.\ de $\overline{\mathbb{H}}$, \uncertain{s'identifie au} groupe des autom.\ de la droite affine $\mathbb{C}$ qui commutent à $\sigma : z \mapsto \bar{z}$, est le \emph{groupe affine réel}.

\page{43}
\note{en haut à gauche, un « 18 » cerclé, de sa main.}
b') Cas $\mathfrak{X} \simeq \mathbb{P}^{1}_{\mathbb{C}}$, $\sigma$ une inversion --- laquelle définit une structure de droite projective réelle --- on trouve \struck{\ill{}} pour $X$ un \emph{disque}. Pour obtenir une métrique distinguée sur $\mathfrak{X}$ invariante par $\sigma$, il faut se donner un pt du disque ouvert $X \setminus \partial X$.

\note{un trait horizontal sépare ce qui suit.}

En résumé, les surfaces conformes à bord connexes admettent une unique métrique riemannienne distinguée, sauf dans les quatre cas suivants
\begin{itemize}
  \item[a) a')] Cas $X \simeq \mathbb{C}$, où $X \simeq \mathbb{H} = \{z \in \mathbb{C} \mid \Im z \geq 0\}$ [droite affine complexe, demi-plan : bord défini par une droite affine réelle], on a alors une classe de métriques mod.\ constantes, qui forment une orbite sous $\operatorname{Aut} X$
  \item[b) b')] $X \simeq \mathbb{P}^{1}_{\mathbb{C}}$ --- $X \simeq$ disque [droite projective complexe \struck{\ill{}} défini par droite projective réelle] --- il y a alors un ens.\ de métriques distinguées dépendant de \add{3} paramètres réels \ill{} \uncertain{des} 2 paramètres réels, également orbite sous $\operatorname{Aut}$.
\end{itemize}

\page{44}
\note{feuillet d'un autre papier ; en haut, au centre, un numéro cerclé qu'on lit « $1'$ » : le signe qui suit le 1, court et levé, ressemble au prime du « $2'$ » de la page 45, et « 17 » n'est pas tout à fait exclu. Par son texte, ce feuillet intercalaire se place après le feuillet 1 de l'auteur (page 26) ; c'est une induction, la page ne le dit pas.}
\textbf{Corollaire 3.} Soit $X \subset X'$ \add{ouvert 0-connexe}, $X'$ hol.\ 0-connexe. Alors $X$ est hyperbolique, sauf dans les cas suivants :
\begin{itemize}
  \item[a)] $X'$ est $\simeq \mathbb{P}^{1}_{\mathbb{C}}$ (resp.\ $\mathbb{E}^{1}_{\mathbb{C}}$, resp.\ $\mathbb{G}_{m\,\mathbb{C}}$), et $\operatorname{card}(X \setminus X') \leq 2$ (resp.\ $\leq 1$, resp.\ $= 0$)\note{sic : $X \setminus X'$ ; le sens demande $X' \setminus X$.} ;
  \item[b)] $X'$ est compacte de genre $1$, et $X = X'$.
\end{itemize}

\textbf{Remarque.} Soit $X$ une surface conforme 0-connexe (pas nécessairement orientable). On dit encore que $X$ est \emph{hyperbolique} si $\widetilde{X} \simeq \mathbb{H}$, ou ce qui revient au même, si toute applic.\ conforme $\mathbb{C} \to X$ est constante.
\marginal{ou encore si le rev.\ \struck{\ill{}} des orientations de $X$ (qui est une surface holomorphe (\uncertain{orient.})) est hyperbolique.}

\struck{\uncertain{Ainsi}} Les surfaces conformes \add{sans bord} 0-connexes \emph{non} hyperboliques sont, en plus des quatre cas énumérés dans (i) ci-dessus, les quotients de telles surfaces par des antiinvolutions sans pts fixes. On trouve donc de plus

\page{45}
\note{même papier que la page 44 ; en haut, au centre, un « $2'$ » cerclé, de sa main. Par son texte, ce feuillet intercalaire se place après le feuillet 2 de l'auteur (page 27) ; c'est une induction, la page ne le dit pas.}
Dans le cas où $X$ est compacte, on trouve qu'il y a \uncertain{la seule} métrique riemannienne (compatible avec la structure conforme) qui soit de courbure constante égale à $-1$.

2) Cas \struck{elliptique} compact connexe (\uncertain{orientable} \ill{}) de genre $1$. Dans le cas orienté, i.e.\ celui d'une courbe elliptique complexe, dont le rev.\ universel est une droite complexe, on a sur celui-ci une métrique canonique mod.\ scalaire (métriques euclidiennes) invariante par translations (et qui \uncertain{l'est} par tout \uncertain{autom.}\ i.e.\ par transformation affine), qui \uncertain{passe} descendre à $X$. La constante multiplicative
\note{la phrase s'interrompt au bas du feuillet ; la suite est vraisemblablement le haut de la page 28 (lot 2), « déterminée par la condition que l'aire totale de $X$ soit $= 1$ ».}

\page{47}
\note{feuillet de listing, sur un autre papier que les précédents ; en haut, de sa main : « 15.1. » (jour et mois, sans année). Le sujet change : ce feuillet porte sur les isomorphismes de Teichmüller et les groupes fondamentaux, non sur les surfaces conformes à bord.}
\begin{itemize}
  \item[] $X$ lisse type $(g,\nu)$ \struck{non singulière} (sur $\overline{\mathbb{Q}}$)
  \item[] $Y$ stable non lisse type $(g,\nu)$ avec propagateur \struck{normal} $N$, d'où un groupe fondamental ext.\ \struck{modifié} $\pi_1([Y,N])$
  \item[] $X \overset{\tau}{\underset{T}{\simeq}} Y$ un isomorphisme de Teichmüller
\end{itemize}

NB Cet isomorphisme induit un isom.\ de groupes extérieurs \struck{\ill{}} (et est déterminé par lui)
\[
  \pi_1(X) \longrightarrow \pi_1([Y,N])
\]
Quand on se donne un « pt base sur $X$ » (ordinaire, ou à l'infini), et un sur $[Y,N]$ (idem --- et de plus distinct des pts singuliers de $Y$) on a une notion d'isom.\ de Teichmüller respect.\ \add{à} les pts marqués, \struck{\ill{}} \add{\uncertain{lesquels}} \uncertain{peuvent} être tous deux à distance finie, ou tous deux à l'infini) et on trouve alors un isom.\ de groupes (pas ext.\,!)
\[
  \pi_1(X,x) \longrightarrow \pi_1([Y,N],y)
\]
En tout état de cause, \uncertain{cet} isom.\ extérieur déjà n'est pas compatible avec les actions ext.\ de $\operatorname{Gal}(\overline{\mathbb{Q}}/K)$, $K$ \add{\uncertain{ext.}}\ \uncertain{assez} extension finie de $\overline{\mathbb{Q}}/\mathbb{Q}$ assez grande.

NB Mettre en évidence \emph{trois} classes de conjugaison d'\struck{\ill{}} hom.\ $\hat{\mathbb{Z}}(1) \to \pi_1(X)$ (\uncertain{deux seulement} si $\nu = 0$, l'une des trois étant les « \uncertain{bouts} »).

\page{48}
\emph{À examiner :}
\begin{itemize}
  \item[a)] Extension de la conjecture fondamentale en géom.\ alg.\ anabélienne aux courbes anabéliennes \struck{stables} singulières (stables, avec propagateur normal). ?
  \item[b)] Pour une $Y$ (singulière anabélienne) donnée (sur corps $K$ de t.f.\ sur $\mathbb{Q}$), \struck{est-il} est-il vrai qu'on récupère les pts de $Y_0 = \operatorname{Sing}(Y) \cup \mathrm{Pts}$ marqués en termes des classes de conj.\ \struck{de} sous-groupes de $\pi_1(Y)$, \uncertain{isom.}\ à $\hat{\mathbb{Z}}$, et \uncertain{c'est} normalisés par $\pi_1(K)$ (l'opération de $\pi_1(K)$ dessus se fait automatiquement via le caractère cyclotomique) ?
  \item[c)] S'inspirer des compactifications MD des $M_{g,\nu}$, définir ce que c'est qu'une compactification anabélienne d'une variété lisse anabélienne (sur un corps $K$ de car.~$0$, quitte peut-être à supposer $K$ de type fini sur $\mathbb{Q}$). Comportement par spécialisation ???
\end{itemize}
\note{« MD » : ainsi sur la page, sans doute pour Mumford--Deligne.}

\page{49}
\note{même papier de listing, écrit sur l'endroit imprimé.}
Générateurs et relations pour les groupoïdes et groupes de Teichmüller (spéciaux et ordinaires, \uncertain{év.}\ toriques…)

Je m'attends à trouver trois types de présentation, à propriétés complémentaires --- mais toutes à signification intrinsèque :

a) Présentation par « \emph{pantalons} » ou « \emph{triades} ». Elle revient à travailler avec les pts à l'infini « les plus antiques » sur les multiplicités modulaires, et semble (d'après l'exemple de $T^{+!}_{0,4}$) commode surtout pour exprimer des groupes du type $T^{+!}_{g,\nu}$ (ou $T^{+}_{g;I}$).

b) Présentation par « boîte » complète de pièces de Légo (voire des boîtes encore plus complètes), à l'aide des \struck{antiautomorphismes} \add{groupes d'}automorphismes de ces pièces, et certains « chemins » de déformation particuliers allant de certaines pièces à d'autres. Exemples frappants : $T^{+}_{0,4}$ (\uncertain{sans}\,!), ou $T^{+}_{1,1} \simeq \mathrm{Sl}(2,\mathbb{Z})$ engendré par $\rho, \sigma$ avec $\rho^{3} = \sigma^{2}$, $\rho^{6}\ (= \sigma^{4}) = 1$.

c) Présentations par \emph{involutions}, qui sont dans $T^{-}_{g,\nu}$. Exemple typique $T_{1,1} \simeq \mathrm{Gl}(2,\mathbb{Z})$ \struck{\ill{}} (\struck{\ill{}} \uncertain{Groupe \emph{antiholomorphique}} triangulé) engendré par trois réflexions $r_1, r_2, r_3$ avec
\[
  (r_1 r_2)^{2} = (r_2 r_3)^{3}, \qquad (r_1 r_2)^{4}\ \bigl(= (r_2 r_3)^{6}\bigr) = 1 .
\]
\note{les exposants de cette dernière ligne sont surchargés (un $4$ sur un $2$, un $3$ sur un autre chiffre) ; la lecture donnée est celle de la dernière encre.}

\page{50}
\note{feuille de brouillon, écrite en tous sens sur l'endroit imprimé d'un listing (qui porte, imprimée, la date du 16 août 1982) : des formules isolées, des suites exactes et trois croquis de sphères au crayon. On les donne dans l'ordre de la feuille retournée, de haut en bas et de gauche à droite ; leur disposition n'est pas reproduite.}

$\mathcal{T}^{+!}_{0,4} \simeq \Pi_{0,3}$ \struck{\ill{}}

\begin{tikzcd}[column sep=small, row sep=small, nodes={font=\scriptsize}]
1 \arrow[r] & \Pi_{0,3} \arrow[r] \arrow[d, no head, "\Vert" description] & \mathcal{T}^{+}_{0,4} \arrow[r] \arrow[d, "\mathrm{can}"] & \mathfrak{S}_4 \arrow[r] \arrow[d, "\mathrm{can}"] & 1 \\
1 \arrow[r] & \Pi_{0,3} \arrow[r] & \mathcal{T}^{+}_{0,1}/\pm 1 \arrow[r] & \mathfrak{S}_3 \arrow[r] & 1
\end{tikzcd}

\note{sous la seconde ligne : $\mathcal{T}^{+}_{0,1}/\pm 1 \simeq \mathrm{Sl}(2,\mathbb{Z})/\pm 1$ et $\mathfrak{S}_3 \simeq \mathrm{Sl}(2,\mathbb{Z}/2\mathbb{Z})$. L'indice « $0,1$ » est lu tel quel ; l'isomorphisme avec $\mathrm{Sl}(2,\mathbb{Z})/\pm 1$ attend plutôt $\mathcal{T}^{+}_{1,1}$, ici comme dans les deux lignes suivantes.}

$\mathcal{T}^{+\prime}_{0,4} \simeq \mathcal{T}^{+}_{0,1}/\pm 1 \simeq \mathrm{Sl}(2,\mathbb{Z})/\pm 1$

$0 \to \Pi_{0,3} \to \mathcal{T}^{+}_{0,1} \to$ \struck{\ill{}} $\widetilde{\mathfrak{S}}_3 \to 1$, avec $\mathcal{T}^{+}_{0,1} = \mathrm{Sl}(2,\mathbb{Z})$ et $\widetilde{\mathfrak{S}}_3 = \mathrm{Sl}(2,\mathbb{Z}/4\mathbb{Z})/\ldots$ \note{sous $\Pi_{0,3}$, une courte ligne biffée, illisible ; à côté, $\mathrm{Gl}(2,\mathbb{F}_2)$ surchargé.}

$\mathrm{Sl}(2,\mathbb{Z})/(\pm 1.\mathbb{Z}) \simeq \mathbb{P}^{1}_{\mathbb{Z}} \simeq \mathbb{P}^{1}(\mathbb{Q})$, le sous-groupe $\pm 1.\mathbb{Z}$ marqué « Borel », et une flèche vers $\mathbb{P}^{1}(\mathbb{Z}/2\mathbb{Z})$.
\[
  \ell_0\,\ell_1\,\ell_2\,\ell_3 = 1, \qquad \ell_0\,\ell_1\,\ell_\infty = 1, \qquad \ell'_0\,\ell'_2\,\ell'_3 = 1, \qquad \ell_\infty = \ell'_0
\]
\note{les indices primés de la troisième relation sont surchargés, et la dernière égalité est lue sous réserve.}

$\mathcal{C}$, \quad $\Pi_{0,3}$, \quad $\mathcal{C}SD$, \quad $\mathcal{C}S$, \quad $\mathcal{C}D$, \quad $\mathbb{T}'SD$, \quad $\mathbb{T}S$, \quad $\mathbb{T}D$, \quad $\mathbb{T}'$, \quad $\mathcal{T}^{+\prime}_{0,4}$

\begin{itemize}
  \item[] groupoïdes : $\mathcal{T}$, $\mathcal{T}^{+}$ \struck{\ill{}}
  \item[] groupes : $\mathbb{T}$ ; $T^{+}$, $T^{!}$, $T^{+!}$, $T'$, $T'^{+}$
  \item[] multiplicités : $\mathbb{M}$
  \item[] \emph{espace de $T$} : $\widetilde{\mathbb{M}}$
\end{itemize}

$\mathrm{Gl}(2,\mathbb{Z})/\mathbb{Z}^{*}$

\[
  1 \to \Pi_{g,\nu-1} \to \mathbb{T}_{g,\nu} \to \mathbb{T}_{g,\nu-1} \to 1
\]
avec, au-dessous, $\Pi_{g,\nu-1}$ rapporté à $X_{g,\nu-1}$, un « $\wr$ » vers $\pi_1(\mathfrak{X}_{g,\nu-1}, \mathbb{T}_{g,\nu-1}; s_\nu)$, et $X_{g,\nu-1} \subset \mathfrak{X}$.

$M_{g,\nu}$ --- \quad $\Pi_{0,3}$ \quad $\mathfrak{S}_3$

\note{les trois croquis au crayon : une sphère portant les points marqués $0$, $1$, $\infty$ et $\alpha$, traversée de grands cercles ; un ovale seul, marqué $\alpha_0$ ; une sphère dont l'équateur porte $0$ et $1$, le pôle $\infty$, reliés par des arcs.}

\page{52}
\note{feuillet de brouillon, très surchargé dans sa moitié supérieure ; on donne la dernière rédaction, les correspondances qu'il écrit par des flèches doubles.}
\begin{itemize}
  \item[] Cercles toriques $\longleftrightarrow$ \struck{\ill{}} \uncertain{plans} vectoriels \add{réels} unidimensionnels \struck{\ill{}} $\longleftrightarrow$ [Droites vectorielles \uncertain{unitaires} complexes $C$ iso.\ à $\mathbb{C}$ (mais pas canon.)]
  \item[] Cercles toriques orientés $\longleftrightarrow$ \struck{\ill{}} \uncertain{plans} vectoriels réels orientés $\longleftrightarrow$ \struck{\ill{}} droites vectorielles complexes unitaires \struck{\ill{}} i.e.\ avec iso $T \simeq T^{\vee}$ i.e.\ $T^{\vee} \simeq \overline{T}$
  \item[] $\updownarrow$
  \item[] torseurs sous $\mathbb{U}$ $\longleftrightarrow$ droites vectorielles complexes \uncertain{conformes}, avec \struck{\ill{}} \struck{\ill{}} contour fixé $D$
\end{itemize}
\marginal{\ill{} \uncertain{toriques} \ill{} \ill{}}

\struck{Le groupe des \ill{} de catégories}

\struck{Soit $C'$, $C''$ deux cercles toriques orientés}, \struck{\ill{}}
\struck{$\mathrm{Isom}(C'^{-1}, C'') \simeq \mathrm{Isom}(C''^{-1}, C')$}
\marginal{\ill{} \ill{} des $\mathbb{R}$}

NB Pour $D$ droite complexe fixée, on a \uncertain{tous les \ill{}} $= \mathbb{D} \times \mathbb{U}$, et $\mathfrak{I}(D) \simeq \mathfrak{I}(D_{\mathbb{U}})$ (\uncertain{avec} $=$ groupe de rotations \uncertain{$\mathbb{U}_{\varepsilon}$}).

Le groupe des cercles toriques \add{orientés} dual (correspond au \uncertain{passage} complexe dual \struck{\ill{}} $T^{\vee}$ \struck{conjuguée} $(\simeq \overline{T})$ i.e.\ l'inverse \uncertain{naturel} \struck{\ill{}}

La donnée, pour deux \emph{droites} vect.\ complexes $T'$, $T''$ d'un système $((D', D''), u)$, où $D'$, $D''$ sont des disques \uncertain{dans} $T'$, $T''$ et $u \in \partial D' \wedge_{\mathbb{U}} \partial D''$
\[
  \bigl(\simeq \mathrm{Isom}_{\mathbb{U}}(\partial D'^{-1}, \partial D'') \simeq \mathrm{Isom}_{\mathbb{U}}(\partial D''^{-1}, \partial D')
\]
\[
  \simeq \mathrm{Isom}(T'^{-1}, T'') \simeq \mathrm{Isom}(T''^{-1}, T') \simeq \mathrm{Bi}\ldots(T' \otimes T'')
\]
\[
  \simeq T'^{\times} \wedge_{\mathbb{C}^{\times}} T''^{\times} \bigr) \simeq
\]
\struck{\ill{}} dualités entre $(T', T'')$)
revient à la donnée d'une $v \in (T' \otimes_{\mathbb{C}} T'')^{*}$ \uncertain{i.e.}\ une dualité entre $T'$, $T''$, \emph{plus} une \uncertain{\ill{}} $D'$ dans $T'$ (\ill{} ce qui revient au même \ill{}, \uncertain{plus} un \uncertain{disque} $D''$ dans $T''$).
\note{« $\mathrm{Bi}\ldots$ » : le mot est abrégé, sans doute \uncertain{Bilin}, des formes bilinéaires ; lecture incertaine.}

\section*{Chirurgie conforme}

\page{53}
\note{en tête du feuillet, de sa main : « 11.1.84 » à gauche, « 1 » au centre, puis le titre souligné « Chirurgie conforme », qu'on a pris pour titre de section. Il numérote en tête les rectos 53, 55, 57 et 59, de 1 à 4 (au 59, le « 4 » suit un chiffre biffé) ; les pages 54, 56 et 58 ne portent pas de numéro, le chiffre de la page 60 est raturé, et la suite reprend au « 5 » de la page 61 (lot 4).}
\uncertain{(1)} \emph{Surface conforme à bord.} On \struck{prendra} la définit par des cartes locales qui, au voisinage d'un point du bord, sont modelées sur $\overline{H} = (\Im z \geq 0) \subset \mathbb{C}$. Dans le cas $X$ orienté i.e.\ structure holomorphe, la structure peut \struck{se} définir par un faisceau de $\mathbb{C}$-algèbres, qui sur le modèle $\overline{H}$ précédent, aux points $x$ de $\partial \overline{H} = \mathbb{R}$, est formé des germes de fonctions qui se \emph{prolongent} en fonctions hol.\ \struck{aux pts de} \add{au} voisinage de $x$ \emph{dans} $\mathbb{C}$. À ne pas confondre avec le sous-faisceau plus grand du faisceau des fonctions complexes continues, savoir \struck{celui} celui des fonctions continues qui sont hol.\ dans $\operatorname{Int}(X)$.
\note{le numéro cerclé, en tête d'alinéa, peut se lire « 4 » ; la suite des numéros cerclés (2), (3), (4) aux p.~55, 57 et 60 demande (1).}

\textbf{Théorème 1.} La structure conforme de $X$ (surface conforme à bord) est connue quand on connaît celle de $\operatorname{Int}(X)$, (ce qui suffit à déterminer le faisceau des fonctions \add{complexes} continues sur $X$ qui sont conformes dans $\operatorname{Int}(X)$).

Il suffit de voir que cela permet de reconstituer les \emph{cartes} conformes des voisinages des points du bord. Je dis que ces cartes

\page{54}
$U \to \overline{H}^{+}$ sont exactement les cartes topologiques (i.e.\ \struck{applications continues injectives} plongements \struck{inj.\ ouverts} ouverts continus inj., \ill{} \struck{bord} du bord) qui sont conformes dans $\operatorname{Int} U$.

Cela \struck{précédent du} équivalant au

\textbf{Lemme 1.} Soit $U$ un ouvert de $\overline{H}^{+}$, \struck{$U \to$ \ill{} injection} plongement ouvert $f : U \hookrightarrow \overline{H}$ \struck{fonction continue sur $V$}, qui est hol.\ dans $U \cap H$, est holo.\ (i.e.\ se prolonge holomorphiquement dans $\mathbb{C}$, au voisinage de $U$ et de $U \cap \partial H = U \cap \mathbb{R}$).

C'est un cas particulier du \emph{lemme de Schwarz}, qui dit que toute fonction \struck{hol.}\ continue complexe dans $U$, qui est \emph{réelle} sur $U \cap \mathbb{R}$, est holomorphe.

Pour terminer la démonstration du \struck{lemme} th., il faut voir encore que la $f$ du lemme est bien (au voisinage de tout pt $x$ de $U \cap \partial H$) une carte holomorphe de $\partial U$ (i.e.\ de $H$) en $x$, i.e.\ que la \uncertain{dérivée} en $x$ de la fonction holomorphe \emph{sur $\mathbb{C}$} définie par $f$ est inversible. Mais cela \ill{} \ill{} par continuité, \struck{et distinction} de ce que \uncertain{son} dévelop-

\page{55}
\note{en tête du feuillet, de sa main : « 2 ».}
\marginal{Notes : \uncertain{indétermination}. \uncertain{Déf.}\ d'une carte loc.\ au bord : modulo \uncertain{chgt.}\ de coordonnées par \uncertain{conf.}\ \ill{} \ill{} \uncertain{réels} de l'autre.}
\add{pe}ment en série qui commence par
\[
  a_n z^{n}, \qquad n > 0,\ a_n \neq 0 \qquad (\text{OPS } x = 0)
\]
on voit tout de suite que l'on doit avoir $n = 1$ ---

[\,\textbf{Thm 2.} Équiv.\ de la cat.\ des surfaces conformes compactes, à bord avec celle des courbes alg.\ propres lisses sur $\mathbb{R}$.\,]

\marginal{NB À partir d'ici, $H$ désigne le \uncertain{demi-plan} fermé $\{z \in \mathbb{C} \mid \Im z \geq 0\}$.}

(2) Découpage d'une surface conforme à bord, suivant un 1-complexe top.\ $K \subset X$, satisfaisant à des hypothèses de régularité convenables (rectifiabilité suffit ?).
\note{le numéro est cerclé, comme le (3) de la p.~57 et le (4) de la p.~60.}

\textbf{Théorème 3.} Le découpé $\widetilde{X}$, muni du \struck{\ill{}} \struck{structure} faisceau des fonctions \add{complexes} continues conformes dans $\operatorname{Int}(\widetilde{X})$ ($\simeq X \setminus K$), est une surface conforme à bord.

\marginal{C'est un th.\ du type « uniformisation \uncertain{locale} » \struck{(à bord)}.}

Idée preuve : vérifier \struck{deux} que pour tout pt $\tilde{x}$ de $\partial \widetilde{X}$, il existe un voisinage ouvert $U$ de $\tilde{x}$, et un \struck{\ill{}} homéo de $U$ avec $H \cap \mathring{D}$ (où $D = \{z \in \mathbb{C} \mid |z| \leq 1\}$ est le disque unité de $\mathbb{C}$). Cela revient à deux choses, suivant que $x$ correspond :

\marginal{\ill{} dans \ill{} Cor.\ au thm.}
\note{en marge du dernier alinéa, un croquis : un demi-disque dont le diamètre porte un point marqué.}

\page{56}
à un pt de $K$ d'ordre $\geq 2$, ou d'ordre $1$ :

\textbf{Lemme 2.} Soit $K$ un segment topologique $\subset \mathbb{C}$ (On suppose $K$ avec régularité (rectifiable, suffit ??)), et $x \in K$. Alors

a) Si $x$ n'est pas extrémité de $K$, alors il existe \struck{un voisinage $U$ de $x$ dans $\mathbb{C}$, ne contenant pas d'extrémités de $U$, tel que $U \setminus K \cap U$ ait deux composantes connexes $U'$ et $U''$ avec $\overline{U'} \cap U = U' \cup (K \cap U)$, $\overline{U''} \cap U = U'' \cup (K \cap U)$, et \ill{} hol.}
\note{ce passage est encadré puis abandonné ; en marge, un croquis : un point $x$ sur un arc, un voisinage en pointillé, une rive hachurée.}
une application continue injective \struck{de $H \cap \mathring{D}$ dans $\mathbb{C}$, hol.\ dans $H \cap$ \ill{}} holomorphe dans $(\operatorname{Int} H) \cap \mathring{D}$, \struck{continue} \struck{$= \,]-1,+1[\, = (D^{++} \cup D^{--})$ \ill{} dans deux \ill{} de $\mathbb{R}$} (mais pas \ill{} hol.\ \ill{}) ($\partial H \cap \mathring{D} = \,]-1,+1[\,$) telle que $f(\partial H \cap \mathring{D}) \subset K$ et $f(0) = x$. De plus, pour une rive donnée de $K$ en $x$, OPS que cette rive est celle déterminée par $\operatorname{Im} f = f(H \cap \mathring{D})$.

b) Si $x$ est extrémité de $K$, \struck{et injective,} alors $\exists$
\[
  f : H \cap \mathring{D} \to \mathbb{C},
\]
hol.\ \struck{dans} $(\operatorname{Int} H) \cap \mathring{D}$, \struck{et $f(-1) = f(1)$ \ill{} injective dans \ill{} telles que} telles que $f(0) = 0$ \struck{\ill{} $f(\,]-1,0)$ \ill{}} \struck{soit \ill{}} $f \mid [-1,0]$ et $f \mid [0,1[$ \uncertain{bijectives}, \ill{}
\note{le bas de la page est une suite de surcharges et de ratures ; la version tenue est au haut de la page suivante. En marge, un croquis : un point $x$ au bout d'un arc, dans un disque.}

\page{57}
\note{en tête du feuillet, de sa main : « 3 ».}
et telles que
\begin{itemize}
  \item[a)] $f \mid (\operatorname{Int} H) \cap \mathring{D}$ est injective, son image $\subset \mathbb{C} \setminus K$ ;
  \item[b)] $f(0) = x$, \struck{$f(-t) = f(t)$} ;
  \item[c)] $f \mid \,]-1,0]$ et $f \mid [0,1[$ injectives, avec même image $\subset K$.
\end{itemize}

\textbf{Scholie.} Chaque rive de $K$ \struck{est} munit $K$ d'une structure analytique réelle canonique, vu l'inclusion de $K$ dans $\mathbb{C}$ (ou dans une surface conforme quelle), mais le passage de l'une à l'autre se fait par une application (identique) qui n'est pas iso analytique : les propriétés de régularité (fonction $C^{\infty}$, ou analytique réel etc.) \uncertain{reflètent} \ill{} les propriétés de régularité correspondantes de $K$ dans $\mathbb{C}$. (Ceci devrait se \ill{} dans \struck{le th.} \uncertain{un} Scholie \uncertain{plus} circonstancié au théorème suivant).

(3) \emph{Recollement.} On doit avoir un théorème dans le style du

\page{58}
\textbf{Th.~4.} Équivalence de catégories entre
\begin{itemize}
  \item[a)] Cat.\ des couples $(X, K, S)$, $X$ surface conforme à bord compacte, $(K, S)$ complexe top.\ \add{compact} \struck{ext} avec
  \[
    \partial X \subset K \subset X,
  \]
  $K$ avec régularité (rectifiable suffit) ;
  \item[b)] Cat des systèmes $(\widetilde{X}, \widetilde{S}, \sigma)$, où $\widetilde{X}$ est une surface conforme à bord \add{compacte}, $\widetilde{S}$ une partie finie de $\partial \widetilde{X}$, $\sigma$ une involution de $\text{Déc}(\partial \widetilde{X} / \widetilde{S})$, $\sigma$ avec régularité (rectifiable suffit).
\end{itemize}

\marginal{NB On ne suppose pas \ill{} les \ill{} des $K \setminus S$ dans \uncertain{les intervalles} b) (il \uncertain{peut y avoir} des cercles).}

Le foncteur a) $\to$ b) résulte du th.~3, avec un complément qui explicite les propriétés de régularité de la fonction-recollement dont il a été question dans le lemme~2. Pour le foncteur en sens inverse, il est connu topologiquement, et comme faisceau structural conforme (\ill{}) sur $X$ ($= \widetilde{X}/\mathcal{R}$), on prend les fonctions continues qui sont conformes dans $X \setminus K \simeq \operatorname{Int} \widetilde{X}$.

\page{59}
\note{en tête du feuillet, de sa main : un chiffre biffé, puis « 4 ».}
Le fait que ça donne bien une structure conforme à bord \struck{équivalent} (qui automatiquement redonne \struck{naissance} \uncertain{avec $K$} à la structure conforme à bord $\widetilde{X}$ de départ) équivaut aux \struck{lemmes suivants} deux cas particuliers suivants :

\struck{\textbf{Lemme 3.} a) Soit $\sigma$ un homéomorphisme rectifiable (ou plus régulier que ça ?) de $[0,1]$ sur lui-même, avec $\sigma(1) = 1$, alors il existe une \struck{courbe} partie fermée $K \subset H \cap \mathring{D}$, \struck{bien contenant $0$, telle que $K \setminus \{0\}$ soit contenue dans $\operatorname{Int} H$} $K \cap \partial H = \{0\}$, $K$ étant un intervalle \ill{} ayant $0$ comme extrémité, qui partage $H \cap \mathring{D}$ en deux composantes $U'$, $U''$. De plus, on trouve un}
\note{ce premier « Lemme 3 » est barré de trois longs traits obliques. En marge, deux croquis : deux demi-plans hachurés dont les bords, marqués d'un point et de l'extrémité $1$, sont reliés par une flèche verticale ; un demi-disque partagé par un arc issu d'un point du diamètre.}

\struck{comme} \textbf{Lemme 3} (en particulier \emph{(deux exemplaires)}),

a) Recollement de $H$ : $H$ suivant un iso.\ rectifiable $[0,1] \xrightarrow{\sigma} [0,1]$, régulier au voisinage de $0$ [ou variante (\struck{\ill{}} utilisée th.~3) : suivant un iso rectifiable $[-1,1] \xrightarrow{\sigma} [-1,1]$, régulier au voisinage de $0$].

\page{60}
\note{en tête du feuillet, de sa main : un chiffre, raturé.}
b) Recollement de $\overline{H}$ : lui-même suivant iso.\ rectifiable $[0,1] \xrightarrow{\sigma}$ \uncertain{$[1, 0]$} ($\sigma(0) = 0$), régulier au voisinage du pt $0$.
\note{les bornes de l'intervalle d'arrivée sont surchargées l'une et l'autre.}

\marginal{12.1 \quad \struck{(a)} c) Soit $C$ une \uncertain{courbe} \ill{} \uncertain{dans} un ouvert $U$ de $\mathbb{C}$, \struck{$f$ un germe de fonction} $f$ une fonction \uncertain{complexe} continue sur $U$, qui est \uncertain{continue} dans $U \setminus C$, alors $f$ est holomorphe dans $U$.}
\note{la note de marge, écrite en oblique, est reliée par un trait au numéro (4) ci-dessous.}

(4) \emph{Recollements \struck{bord à bord} de disques.}
\marginal{Le recollement « bord à bord ».}
Cela correspond, dans le th.~4, au cas où $K$ est une variété \add{(compacte)} de dim~$1$, $S = \emptyset$. \struck{Il} Pour avoir « des cas \ill{} \ill{} », il est naturel de dégager des conditions restrictives naturelles sur l'involution $\sigma$ de $\partial \widetilde{X}$.

Pour ceci, \struck{\ill{}} on se rappelle que pour une surface conforme à bord $X$ (\ill{} de $\widetilde{X}$) et une comp.\ connexe du bord, soit $C$, il revient au même de se donner une structure de droite projective (rectifiable) sur $C$, ou un « \uncertain{bouchage} de trou » sur $C$ \add{de bord $C$} par un disque ; les propriétés de régularité (p.\,ex.\ : la structure analytique réelle de $C$) de la structure de droite projective, reflètent

\end{document}
